% !TeX root = ../教材题库_全部条目_转写工作稿.tex
% Source: XJ-B1-C06-SU-P262-Y1504
% 第6章 / 小结与复习 小结与复习 / 教材 p.262、263 / PDF 第269、270页
% OCR未校对：公式与文字均可能有误；未完成数学TeX转写；禁止直接发给学生
% 原题及相邻原解截图（核对用）：../题目原图/XJ-B1-C06-SU-P262-Y1504-01.webp, ../题目原图/XJ-B1-C06-SU-P262-Y1504-02.webp
% 本文件只在转写完成后作为题干使用；OCR 条目尚不是可靠题干。
\begin{bankitem}{第6章 小结与复习 第3题}{第6章 / 小结与复习 小结与复习 / 教材 p.262、263 / PDF 第269、270页}{XJ-B1-C06-SU-P262-Y1504}{OCR未校对：公式与文字均可能有误；未完成数学TeX转写；禁止直接发给学生}
\begin{ocrdraft}
[教材p.262 / PDF第269页]
\par
3．获取的数据尽管杂乱无章，但经过统计图表的科学呈现，显得井然有序并
\par
体现出诸多有价值的信息．尝试总结本章所介绍的各种统计图表的差异．
\par
[教材p.263 / PDF第270页]
\end{ocrdraft}
\end{bankitem}
